\begin{center}
  {\zihao{2}\bfseries 基于场景随机规划与滚动优化的微网购电调控策略\par}
  \vspace{1.2em}
  {\zihao{3}\bfseries 摘要}
\end{center}

本文研究含光伏与储能的微网外购电计划问题。针对负荷、光伏预测和电价信息逐步丰富的四个任务，建立一套以10 min为步长的统一母线能量平衡模型，并用确定性混合整数规划、两阶段随机规划和随机模型预测控制依次扩展。模型严格区分功率与时段电量、144个流量时段与145个库存状态，并以信息发布时间约束预测和决策，避免使用未来实际数据。

针对问题一，将全天负荷、光伏预测与分时电价视为已知，以购电费最小为目标建立储能调度MILP。计算得到全天购电费为35126.95元，较不使用储能的可行基线降低26.898\%；计划购电量为59482.70 kWh，期初和期末储能量均为6000 kWh。LP与MILP费用一致，最终采用具有严格充放电互斥约束的MILP结果。

针对问题二，在每日0:00仅使用此前历史，比较加权历史均值、傅里叶--ARIMA和随机森林三类预测器，并以1月21--31日滚动费用固定选择随机森林及1.2倍终端价值。最近28个历史日的负荷--光伏成对样本外残差按整日联合抽样256次，形成保持日内及跨变量相关性的场景。以购电和名义电池轨迹为第一阶段、紧急购电与富余处置为第二阶段，连续回放2月至12月334天。主方案实际购电费为1447.93万元，紧急购电量为262.35 MWh；相对同一预测器的点预测方案费用降低7.69\%，相对同初始库存的旧机制对照降低13.653\%。

针对问题三，构造0:00、6:00、12:00、18:00四次更新的随机滚动模型，以历史样本外误差校准不同提前量的偏差与尺度，并用联合日块生成剩余时域场景。风险中性的条件分箱尺度主方案费用为1662.77万元、紧急购电费为91.25万元；四次预报相对仅0:00新预报累计节省41.21万元。光伏正值样本的CRPS为184.54 kW，90\%经验区间覆盖率为80.67\%，显示概率区间仍偏窄。CVaR方案能降低紧急费用及其上尾，但增加平均现金费用。针对问题四，建立价格--净负荷联合场景，并分别提出经历史费用校准的价量依赖收缩模型和只前瞻下一次信息更新的条件调整模型；由于尚未完成求解，本文保留结果表和结论占位，不报告未经计算支撑的节省率。

\vspace{0.8em}
\noindent\textbf{关键词：}微网调度；储能；两阶段随机规划；滚动时域优化；预测不确定性；CVaR
